% --- Archivo: 13_sqp_teoria_basica.tex ---

\section{Programación cuadrática secuencial}\label{v2:sec:4.6}

Como sugiere el nombre, los métodos de programación cuadrática secuencial (SQP)\footnote{En inglés: SQP - sequential quadratic programming} consisten en la resolución de una sucesión de problemas de programación cuadrática (i.e., de minimización de una función cuadrática sujeta a restricciones lineales), que aproximan, en cada iteración, el problema original dado. Un problema de programación cuadrática adecuadamente construido puede ser una buena aproximación del problema original, al mismo tiempo siendo de resolución relativamente fácil. En particular, problemas de programación cuadrática admiten varios métodos bien eficientes, inclusive algunos con convergencia finita (vea el \S 6.2). Por eso, esta estrategia tiene sentido práctico, diferentemente de la utilización de aproximaciones de tercer orden, cuarta, etc.

Por otro lado, métodos SQP pueden ser pensados como una generalización natural de la idea fundamental del método de Newton clásico. A propósito, como veremos más adelante, en el caso de restricciones de igualdad, los métodos SQP pueden ser considerados como una implementación específica del método de Newton aplicado al sistema de Lagrange asociado al problema. En el caso en que hay restricciones de desigualdad, esta equivalencia deja de existir, mas la relación con el método de Newton es clara en este caso también.
Los métodos SQP están considerados entre los más eficientes métodos de optimización de uso general. Para resolver un problema dado, sin estructura especial, con frecuencia son escogidos métodos de esta clase.

\subsection{Restricciones de igualdad}\label{v2:sec:4.6.1}

Consideremos el problema
\begin{equation}
    \min f(x) \quad \text{sujeto a} \quad x \in D = \{x \in \R^n \mid h(x) = 0\},
    \label{v2:eq:4.221}
\end{equation}
donde $f \colon \R^n \to \R$ e $h \colon \R^n \to \R^l$ son funciones dos veces diferenciables en $\R^n$.
Sea $x^k \in \R^n$ una aproximación del punto estacionario $\bar{x}$ del problema \eqref{v2:eq:4.221}. Vamos a aproximar el problema original, en torno de $x^k$, por un problema de programación cuadrática de la forma
\begin{equation}
    \min f(x^k) + \langle f'(x^k), x - x^k \rangle + \frac{1}{2}\langle H_k(x - x^k), x - x^k \rangle
    \label{v2:eq:4.222}
\end{equation}
\[
    \text{sujeto a } x \in D_k = \{x \in \R^n \mid h(x^k) + h'(x^k)(x - x^k) = 0\},
\]
donde $H_k \in \R(n, n)$ es una matriz simétrica.

La primera tentativa, que parece natural, sería utilizar la matriz $H_k = f''(x^k)$, i.e., tomar como función objetivo del subproblema la aproximación de segundo orden de la función objetivo $f$ del problema original (comparar con el método de Newton restringido en la \cref{v2:sec:4.1.3}). Sin embargo, consideraciones más finas muestran que esta elección, aparentemente natural, no es de las mejores. Esto tiene a ver con la siguiente observación. La aproximación lineal de las restricciones no puede proporcionar información sobre la curvatura de las mismas. No obstante, intuitivamente, es claro que esta información de segundo orden sobre las restricciones es tan importante como la información de segundo orden sobre la función objetivo (son dos partes de problema distintas, una tan importante cuanto a otra). Agregar términos cuadráticos a las restricciones de los subproblemas va a tornarlos bien más difíciles de resolver. Actualmente existen algunas propuestas de métodos de este tipo, que hasta tienen ciertas ventajas teóricas. Mas por ahora no puede ser dicho que subproblemas con restricciones cuadráticas se establecieron como una alternativa práctica a los subproblemas con restricciones lineales de SQP. Felizmente, la información de segundo orden sobre las restricciones puede ser repasada al subproblema de otra manera, bastante natural desde el punto de vista Newtoniano, como veremos a seguir. En particular, la propuesta es escoger
\begin{equation}
    H_k = L''_{xx}(x^k, \lambda^k),
    \label{v2:eq:4.223}
\end{equation}
donde $\lambda^k \in \R^l$ es una aproximación al multiplicador de Lagrange $\bar{\lambda}$ asociado al punto estacionario $\bar{x}$, y
\[
    L \colon \R^n \times \R^l \to \R, \quad L(x, \lambda) = f(x) + \langle \lambda, h(x) \rangle,
\]
es la Lagrangiana del problema \eqref{v2:eq:4.221}. O sea, la información de segundo orden sobre las restricciones es pasada al subproblema vía la función objetivo de este. Notamos que el subproblema así obtenido continua a ser de programación cuadrática (con función objetivo dada por la aproximación cuadrática de la función $L(\cdot, \lambda^k)$ en torno de $x^k$, y con el conjunto factible dado por la aproximación lineal de las restricciones en torno de $x^k$).

Otra motivación para esta elección viene dada por la interpretación del problema \eqref{v2:eq:4.222} como una implementación de la iteración de Newton para el sistema de Lagrange del problema original \eqref{v2:eq:4.221}, que fue considerada en el \cref{v2:sec:4.4.1}. Recordamos que el sistema de Lagrange para el problema en cuestión es
\begin{equation}
    L'(x, \lambda) = 0,
    \label{v2:eq:4.224}
\end{equation}
y una iteración de Newton para este sistema consiste en la resolución del sistema lineal
\begin{align}
    L''_{xx}(x^k, \lambda^k)(x - x^k) + (h'(x^k))^\top(\lambda - \lambda^k) &= -f'(x^k) - (h'(x^k))^\top \lambda^k, \label{v2:eq:4.225} \\
    h'(x^k)(x - x^k) &= -h(x^k) \label{v2:eq:4.226}
\end{align}
en relación a $(x, \lambda) \in \R^n \times \R^l$, donde $(x^k, \lambda^k) \in \R^n \times \R^l$ es una dada aproximación a la solución de \eqref{v2:eq:4.224}.
Sea $x^{k+1} \in \R^n$ un punto estacionario del problema \eqref{v2:eq:4.222} e $y^{k+1} \in \R^l$ un multiplicador de Lagrange asociado. Esto significa que $(x^{k+1}, y^{k+1})$ es una solución del sistema
\begin{align*}
    f'(x^k) + H_k(x - x^k) + (h'(x^k))^\top y &= 0, \\
    h(x^k) + h'(x^k)(x - x^k) &= 0
\end{align*}
en relación a $(x, y) \in \R^n \times \R^l$. Comparando este último sistema con \eqref{v2:eq:4.225}-\eqref{v2:eq:4.226}, vemos que ellos coinciden, si la matriz $H_k$ fuese escogida de acuerdo con \eqref{v2:eq:4.223}. Por tanto, esta es la elección que corresponde a la idea de métodos Newtonianos. Podemos, entonces, esperar convergencia local rápida bajo hipótesis naturales.
Formulamos ahora el método de programación cuadrática secuencial (SQP) para problemas con restricciones de igualdad.

\noindent \textbf{Algoritmo 4.6.1. (El método de programación cuadrática secuencial)} \\
Elegir $(x^0, \lambda^0) \in \R^n \times \R^l$ y tomar $k := 0$.
\begin{enumerate}
    \item Calcular $x^{k+1} \in \R^n$, un punto estacionario del problema \eqref{v2:eq:4.222}, donde $H_k$ es dada por \eqref{v2:eq:4.223}, y calcular $y^{k+1} \in \R^l$, un multiplicador de Lagrange asociado a $x^{k+1}$.
    \item Tomar $\lambda^{k+1} = y^{k+1}$.
    \item Tomar $k := k + 1$ y retornar al Paso 1.
\end{enumerate}

Por lo expuesto arriba, para un punto inicial $(x^0, \lambda^0)$ cualquier, el Algoritmo 4.6.1 genera la misma sucesión $\{(x^k, \lambda^k)\}$ que el método de Newton puro aplicado al sistema de Lagrange \eqref{v2:eq:4.224}, ya considerado en el \cref{v2:sec:4.4.1}. Por eso, no hay necesidad de un nuevo análisis (local) del método SQP en este caso. En particular, si las derivadas segundas de las funciones $f$ y $h$ son continuas en el punto $\bar{x}$, y son satisfechas la condición de regularidad de las restricciones \eqref{v1:eq:1.9} y la condición suficiente de segunda orden del Teorema 1.2.2(b), entonces por el Teorema 4.4.1, inmediatamente tenemos que el método SQP converge localmente a $(\bar{x}, \bar{\lambda})$ con tasa superlineal. Más aún, si las derivadas segundas de $f$ y de $h$ son Lipschitz-continuas en una vecindad de $\bar{x}$, entonces la tasa de convergencia es cuadrática.

Como implementación del método de Newton para el sistema de Lagrange, el método SQP tal vez no sea muy útil (por lo menos en el contexto de un estudio local). Esta interpretación es análoga a aquella en el \cref{v2:sec:3.2.2} en el caso del método de Newton para optimización irrestricta y su interpretación vía minimización de funciones cuadráticas. Sin embargo, estas consideraciones son importantes para sugerir estrategias de globalización del esquema local. Además de eso, el método SQP para problemas con restricciones de igualdad sugiere una generalización muy natural para el caso general, cuando hay también restricciones de desigualdad. Esta generalización será presentada a seguir.

% [Ejercicios 4.6.1 y 4.6.2 movidos a exercises.tex]

\subsection{Restricciones mixtas}\label{v2:sec:4.6.2}

Consideremos ahora un problema de optimización con restricciones de igualdad y desigualdad
\begin{equation}
    \min f(x) \quad \text{sujeto a} \quad x \in D = \left\{ x \in \R^n \;\middle|\; \begin{array}{l} h(x) = 0, \\ g(x) \le 0. \end{array} \right\},
    \label{v2:eq:4.227}
\end{equation}
donde $h \colon \R^n \to \R^l$ y $g \colon \R^n \to \R^m$ son dos veces diferenciables en $\R^n$.
Sea $x^k \in \R^n$ una aproximación de un punto estacionario $\bar{x}$ del problema \eqref{v2:eq:4.227}. Una generalización natural del subproblema \eqref{v2:eq:4.222} para el caso de restricciones mixtas es dada por el problema
\begin{equation}
    \min \langle f'(x^k), d \rangle + \frac{1}{2}\langle H_k d, d \rangle \quad \text{sujeto a} \quad d \in D_k,
    \label{v2:eq:4.228}
\end{equation}
\begin{equation}
    D_k = \left\{ d \in \R^n \;\middle|\; \begin{array}{l} h(x^k) + h'(x^k)d = 0, \\ g(x^k) + g'(x^k)d \le 0. \end{array} \right\}.
    \label{v2:eq:4.229}
\end{equation}
El método SQP, por lo tanto, consiste en lo siguiente.

\noindent \textbf{Algoritmo 4.6.2. (El método de programación cuadrática secuencial)} \\
Elegir $(x^0, \lambda^0, \mu^0) \in \R^n \times \R^l \times \R^m_+$ y tomar $k := 0$.
\begin{enumerate}
    \item Calcular un punto estacionario $d^k \in \R^n$ del problema \eqref{v2:eq:4.228}-\eqref{v2:eq:4.229}, donde $H_k \in \R(n, n)$ es una matriz simétrica, y multiplicadores de Lagrange $y^k \in \R^l$, $z^k \in \R^m_+$ asociados a $d^k$.
    \item Tomar $x^{k+1} = x^k + d^k$, $\lambda^{k+1} = y^k$, $\mu^{k+1} = z^k$.
    \item Tomar $k := k + 1$ y retornar al Paso 1.
\end{enumerate}

Por lo observado arriba en el caso de restricciones de igualdad, podemos esperar convergencia rápida del método si escogemos
\begin{equation}
    H_k = L''_{xx}(x^k, \lambda^k, \mu^k),
    \label{v2:eq:4.230}
\end{equation}
donde
\[
    L \colon \R^n \times \R^l \times \R^m \to \R,
\]
\[
    L(x, \lambda, \mu) = f(x) + \langle \lambda, h(x) \rangle + \langle \mu, g(x) \rangle
\]
es la Lagrangiana del problema \eqref{v2:eq:4.227}. Notemos que si $D_k \neq \emptyset$ y la matriz $H_k$ es definida positiva, el subproblema \eqref{v2:eq:4.228}-\eqref{v2:eq:4.229} posee una única solución y esta solución también es el único punto estacionario del problema (por el hecho de que la función objetivo es fuertemente convexa y el conjunto factible es convexo). No obstante, cabe mencionar que bajo hipótesis naturales (tales como la condición suficiente de segunda orden para el problema \eqref{v2:eq:4.227}) no es posible garantizar que la matriz $H_k$ dada por \eqref{v2:eq:4.230} sea definida positiva, mismo para $x^k$ próximo a $\bar{x}$, y $\lambda^k, \mu^k$ próximos a $\bar{\lambda} \in \bar{\mu}$, respectivamente (comparar con la motivación para introducción de Lagrangianas aumentadas en la \cref{v2:sec:4.4.3}, y con el Ejercicio 4.6.1). Por eso, la existencia de puntos estacionarios del subproblema \eqref{v2:eq:4.228}-\eqref{v2:eq:4.229} no es automática y tiene que ser probada como parte del análisis local del método. Además de eso, mismo cuando un punto estacionario existe, él puede no ser único. Por eso, en el análisis local suponemos que $d^k$ sea el punto estacionario del subproblema \eqref{v2:eq:4.228}-\eqref{v2:eq:4.229} de menor norma. En la práctica, este requerimiento es simplemente ignorado.

\begin{theorem}[Convergencia local del método SQP]\label{v2:thm:4.6.1}
    Sean funciones $f \colon \R^n \to \R$, $h \colon \R^n \to \R^l$ y $g \colon \R^n \to \R^m$ dos veces diferenciables en una vecindad del punto $\bar{x} \in \R^n$, con derivadas segundas continuas en este punto. Sea $\bar{x}$ un punto estacionario del problema \eqref{v2:eq:4.227} que satisface la condición de independencia lineal de los gradientes de las restricciones activas \eqref{v1:eq:1.17}.
    Supongamos que $\bar{x}$, con los (únicos) multiplicadores de Lagrange asociados $\bar{\lambda} \in \R^l$ y $\bar{\mu} \in \R^m_+$, satisface la condición suficiente de segunda orden para el problema \eqref{v2:eq:4.227}, dada en el Teorema 1.2.3(b).
    Supongamos, finalmente, la condición de complementaridad estricta
    \begin{equation}
        \bar{\mu}_i > 0 \quad \forall i \in I(\bar{x}) = \{i = 1, \dots, m \mid g_i(\bar{x}) = 0\}.
        \label{v2:eq:4.231}
    \end{equation}
    Entonces para cualquier punto inicial $(x^0, \lambda^0, \mu^0) \in \R^n \times \R^l \times \R^m$ suficientemente próximo a $(\bar{x}, \bar{\lambda}, \bar{\mu})$, el Algoritmo 4.6.2, donde $H_k$ es dada por \eqref{v2:eq:4.230} y $d^k$ es el punto estacionario de \eqref{v2:eq:4.228}-\eqref{v2:eq:4.229} con menor norma, está bien definido. La sucesión generada por el método converge a $(\bar{x}, \bar{\lambda}, \bar{\mu})$. La tasa de convergencia es superlineal. Más aún, si las derivadas segundas de $f, h$ y $g$ son Lipschitz-continuas en una vecindad del punto $\bar{x}$, entonces la convergencia es cuadrática.
\end{theorem}
\begin{proof}
    Denotamos $U = \R^n \times \R^l \times \R^m$. Primero, tenemos que probar que para $(x^k, \lambda^k, \mu^k) \in U$ suficientemente próximo a $(\bar{x}, \bar{\lambda}, \bar{\mu})$, el subproblema \eqref{v2:eq:4.228}-\eqref{v2:eq:4.229} posee un único punto estacionario $d^k$ de menor norma y este punto tiene un único par de multiplicadores de Lagrange $(y^k, z^k)$ asociado. Consideremos el sistema KKT del subproblema \eqref{v2:eq:4.228}-\eqref{v2:eq:4.229}:
    \begin{align}
        f'(x^k) + L''_{xx}(x^k, \lambda^k, \mu^k)d + (h'(x^k))^\top y + (g'(x^k))^\top z &= 0, \label{v2:eq:4.232} \\
        h(x^k) + h'(x^k)d &= 0, \label{v2:eq:4.233} \\
        g(x^k) + g'(x^k)d &\le 0, \quad z \ge 0, \label{v2:eq:4.234} \\
        z_i(g_i(x^k) + \langle g'_i(x^k), d \rangle) &= 0, \quad i = 1, \dots, m, \label{v2:eq:4.235}
    \end{align}
    en relación a $(d, y, z) \in U$. Notemos primero que si $x^k$ es suficientemente próximo a $\bar{x}$ y $d \in \R^n$ tiene norma suficientemente pequeña, no puede existir más de un par de multiplicadores $(y, z) \in \R^l \times \R^m$ tales que el punto $(d, y, z)$ satisfaga \eqref{v2:eq:4.232}-\eqref{v2:eq:4.235}. Con efecto, para tales $x^k$ y $d$, se tiene que
    \[
        g_i(x^k) + \langle g'_i(x^k), d \rangle < 0 \quad \forall i \in \{1, \dots, m\} \setminus I(\bar{x}),
    \]
    por la continuidad de las funciones involucradas, y de \eqref{v2:eq:4.235},
    \[
        z_i = 0 \quad \forall i \in \{1, \dots, m\} \setminus I(\bar{x}).
    \]
    En particular, para el conjunto de los índices $J^k$ de las restricciones activas en la solución del subproblema \eqref{v2:eq:4.228}-\eqref{v2:eq:4.229}, se tiene que $J^k \subset I(\bar{x})$. Los gradientes de estas restricciones forman, entonces, un subconjunto de
    \[
        \{h'_i(x^k), \; i = 1, \dots, l\} \cup \{g'_i(x^k), \; i \in I(\bar{x})\}
    \]
    que es linealmente independiente para $x^k$ próximo a $\bar{x}$ (pela continuidad y por la hipótesis de independencia lineal de estos gradientes en el punto $\bar{x}$). Concluimos entonces que los gradientes de las restricciones activas son linealmente independientes en el subproblema \eqref{v2:eq:4.228}-\eqref{v2:eq:4.229}. Esto implica que no puede haber más de un par de multiplicadores de Lagrange en el sistema \eqref{v2:eq:4.232}-\eqref{v2:eq:4.235}.

    Definimos la función
    \[
        \Phi \colon U \times U \to U,
    \]
    \[
        \Phi(w, u) = \begin{pmatrix} f'(x) + L''_{xx}(x, \lambda, \mu)d + (h'(x))^\top y + (g'(x))^\top z \\
        h(x) + h'(x)d \\
        z_1(g_1(x) + \langle g'_1(x), d \rangle) \\
        \vdots \\
        z_m(g_m(x) + \langle g'_m(x), d \rangle) \end{pmatrix},
    \]
    \[
        w = (x, \lambda, \mu), \quad u = (d, y, z).
    \]
    El sistema
    \[
        \Phi(w, u) = 0,
    \]
    donde $w \in U$ hace el papel de parámetro, corresponde a las ecuaciones \eqref{v2:eq:4.232}, \eqref{v2:eq:4.233}, \eqref{v2:eq:4.235}.
    Por las definiciones de punto estacionario del problema \eqref{v2:eq:4.227} y de multiplicadores de Lagrange asociados, tenemos $\Phi(\bar{w}, \bar{u}) = 0$, donde $\bar{w} = (\bar{x}, \bar{\lambda}, \bar{\mu})$, $\bar{u} = (0, \bar{\lambda}, \bar{\mu})$. Además de eso,
    \[
        \Phi'_u(\bar{w}, \bar{u}) = \begin{pmatrix} L''_{xx}(\bar{x}, \bar{\lambda}, \bar{\mu}) & (h'(\bar{x}))^\top & (g'(\bar{x}))^\top \\
        h'(\bar{x}) & 0 & 0 \\
        \bar{\mu}_1 g'_1(\bar{x}) & 0 & 0 \\
        \bar{\mu}_2 g'_2(\bar{x}) & 0 & 0 \\
        \vdots & \vdots & \vdots \\
        \bar{\mu}_m g'_m(\bar{x}) & 0 & 0 \end{pmatrix}.
    \]
    Tomamos $\tilde{u} = (\tilde{d}, \tilde{y}, \tilde{z}) \in \ker \Phi'_u(\bar{w}, \bar{u})$, arbitrario, i.e.,
    \begin{align}
        L''_{xx}(\bar{x}, \bar{\lambda}, \bar{\mu})\tilde{d} + (h'(\bar{x}))^\top \tilde{y} + (g'(\bar{x}))^\top \tilde{z} &= 0, \label{v2:eq:4.236} \\
        h'(\bar{x})\tilde{d} &= 0, \label{v2:eq:4.237} \\
        \bar{\mu}_i \langle g'_i(\bar{x}), \tilde{d} \rangle + g_i(\bar{x})\tilde{z}_i &= 0, \quad i = 1, \dots, m. \label{v2:eq:4.238}
    \end{align}
    De \eqref{v2:eq:4.238} y de la condición de complementaridad estricta \eqref{v2:eq:4.231}, tenemos que
    \begin{equation}
        \tilde{z}_i = 0 \quad \forall i \in \{1, \dots, m\} \setminus I(\bar{x}), \quad \langle g'_i(\bar{x}), \tilde{d} \rangle = 0 \quad \forall i \in I(\bar{x}).
        \label{v2:eq:4.239}
    \end{equation}
    Además de esto, bajo la condición de complementaridad estricta \eqref{v2:eq:4.231} tenemos que el cono crítico del problema \eqref{v2:eq:4.227} en el punto $\bar{x}$ tiene la siguiente forma:
    \begin{equation}
        \mathcal{K}(\bar{x}) = \{v \in \ker h'(\bar{x}) \mid \langle g'_i(\bar{x}), v \rangle = 0 \quad \forall i \in I(\bar{x})\}.
        \label{v2:eq:4.240}
    \end{equation}
    Por tanto, de \eqref{v2:eq:4.237} y \eqref{v2:eq:4.239}, tenemos que $\tilde{d} \in \mathcal{K}(\bar{x})$. Multiplicando los dos lados de \eqref{v2:eq:4.236} por $\tilde{d}$, obtenemos que
    \begin{align*}
        0 &= \langle L''_{xx}(\bar{x}, \bar{\lambda}, \bar{\mu})\tilde{d}, \tilde{d} \rangle + \langle \tilde{y}, h'(\bar{x})\tilde{d} \rangle + \langle \tilde{z}, g'(\bar{x})\tilde{d} \rangle \\
        &= \langle L''_{xx}(\bar{x}, \bar{\lambda}, \bar{\mu})\tilde{d}, \tilde{d} \rangle + \sum_{i=1}^m \tilde{z}_i \langle g'_i(\bar{x}), \tilde{d} \rangle \\
        &= \langle L''_{xx}(\bar{x}, \bar{\lambda}, \bar{\mu})\tilde{d}, \tilde{d} \rangle,
    \end{align*}
    donde utilizamos \eqref{v2:eq:4.237} y \eqref{v2:eq:4.239}.
    Esto significa que $\tilde{d} = 0$ (pues en caso contrario, la condición suficiente de segunda orden en el punto $\bar{x}$ sería violada). De \eqref{v2:eq:4.236} y \eqref{v2:eq:4.239}, ahora sigue que
    \[
        (h'(\bar{x}))^\top \tilde{y} + \sum_{i \in I(\bar{x})} \tilde{z}_i g'_i(\bar{x}) = 0.
    \]
    Pela condición \eqref{v1:eq:1.17} de independencia lineal de los gradientes de las restricciones activas en $\bar{x}$, obtenemos que $\tilde{y} = 0$ e $\tilde{z}_i = 0 \; \forall i \in I(\bar{x})$. Acabamos de probar que $\tilde{u} = (\tilde{d}, \tilde{y}, \tilde{z}) = 0$, i.e., la matriz $\Phi'_u(\bar{w}, \bar{u})$ es no-singular.
    Aplicando a la función $\Phi$ en el punto $(\bar{w}, \bar{u})$ el Teorema de la Función Implícita (Teorema 1.5.3), concluimos que existen vecindades $\mathcal{W}$ del punto $\bar{w}$ y $\mathcal{V}$ del punto $\bar{u}$ en $U$, para las cuales existe una única función $\chi(\cdot) = (d(\cdot), y(\cdot), z(\cdot)) \colon \mathcal{W} \to \mathcal{V}$ continua en el punto $\bar{w}$ tal que $\chi(\mathcal{W}) \subset \mathcal{V}$ y
    \begin{equation}
        \Phi(w, \chi(w)) = 0 \quad \forall w \in \mathcal{W}.
        \label{v2:eq:4.241}
    \end{equation}
    En particular, $d(\bar{w}) = 0$, $y(\bar{w}) = \bar{\lambda}$, $z(\bar{w}) = \bar{\mu}$. De aquí, de la continuidad de $\chi(\cdot)$ en $\bar{w}$, y de la condición de complementaridad estricta \eqref{v2:eq:4.231}, obtenemos que para una vecindad $\mathcal{W}$ suficientemente pequeña, $\forall w = (x, \lambda, \mu) \in \mathcal{W}$ se tiene que
    \begin{equation}
        g_i(x) + \langle g'_i(x), d(w) \rangle < 0 \quad \forall i \in \{1, \dots, m\} \setminus I(\bar{x}),
        \label{v2:eq:4.242}
    \end{equation}
    \begin{equation}
        z_i(w) > 0 \quad \forall i \in I(\bar{x}).
        \label{v2:eq:4.243}
    \end{equation}
    Llevando en cuenta la definición de la función $\Phi$, las relaciones \eqref{v2:eq:4.241}-\eqref{v2:eq:4.243} significan que, para $w^k = (x^k, \lambda^k, \mu^k) \in \mathcal{W}$, el sistema \eqref{v2:eq:4.232}-\eqref{v2:eq:4.235} posee en $\mathcal{V}$ una única solución, que es $(d(w^k), y(w^k), z(w^k))$. Por el hecho de que los multiplicadores tienen que ser únicos (ya justificado arriba), podemos afirmar que si la vecindad $\mathcal{W}$ es suficientemente pequeña, esta solución es única no solamente en $\mathcal{V}$ mas también en el conjunto $\tilde{\mathcal{V}} \times \R^l \times \R^m$, donde $\tilde{\mathcal{V}}$ es una vecindad de $0$ en $\R^n$. De aquí sigue que el $d^k = d(w^k)$ es el único punto estacionario del problema \eqref{v2:eq:4.228}-\eqref{v2:eq:4.229} de menor norma, y $(y^k, z^k) = (y(w^k), z(w^k))$ es el único par de multiplicadores de Lagrange asociado a $d^k$.
    Observemos que, por \eqref{v2:eq:4.242} y \eqref{v2:eq:4.243}, para $(x^k, \lambda^k, \mu^k) \in \mathcal{W}$ las condiciones \eqref{v2:eq:4.234} y \eqref{v2:eq:4.235} en el sistema \eqref{v2:eq:4.232}-\eqref{v2:eq:4.235}, que define $(d^k, y^k, z^k)$, pueden ser sustituidas por
    \begin{equation}
        g_i(x^k) + \langle g'_i(x^k), d \rangle = 0 \quad i \in I(\bar{x}),
        \label{v2:eq:4.244}
    \end{equation}
    \begin{equation}
        z_i = 0 \quad i \in \{1, \dots, m\} \setminus I(\bar{x}).
        \label{v2:eq:4.245}
    \end{equation}
    En particular, una iteración del método es equivalente al cálculo de la solución $(d^k, y^k, z^k)$ del sistema de ecuaciones lineales \eqref{v2:eq:4.232}, \eqref{v2:eq:4.233}, \eqref{v2:eq:4.244}, \eqref{v2:eq:4.245}.
    Consideremos el problema con restricciones de igualdad
    \[
        \min f(x) \quad \text{sujeto a} \quad x \in \tilde{D},
    \]
    \[
        \tilde{D} = \{x \in \R^n \mid h(x) = 0, \; g_i(x) = 0, \; i \in I(\bar{x})\}.
    \]
    Como es fácil ver, bajo nuestras hipótesis el punto $\bar{x}$ es una solución local de este problema que satisface la condición de regularidad de las restricciones y la condición suficiente de segunda orden para problemas con restricciones de igualdad; vea \S 1.2. Por tanto, si le agregamos al sistema de Lagrange correspondiente las ecuaciones
    \[
        \mu_i = 0, \quad i \in \{1, \dots, m\} \setminus I(\bar{x}),
    \]
    el sistema de ecuaciones en relación a $(x, \lambda, \mu) \in U$ así obtenido tendrá una solución $(\bar{x}, \bar{\lambda}, \bar{\mu})$ que es regular (en el sentido de que la derivada de este sistema en esta solución es no-singular). De acuerdo con el Teorema 3.2.1, el método de Newton para este sistema converge (localmente) a la solución con tasa superlineal, y si las derivadas segundas de $f, h$ y $g$ son Lipschitz-continuas en una vecindad del punto $\bar{x}$, la tasa de convergencia es cuadrática.
    Como es fácil ver, para una aproximación $(x^k, \lambda^k, \mu^k)$ dada, una iteración del método de Newton consiste en un paso al punto $(x^k + d^k, y^k, z^k)$, donde $(d^k, y^k, z^k)$ es la solución del sistema que consiste en las ecuaciones \eqref{v2:eq:4.233}, \eqref{v2:eq:4.244}, \eqref{v2:eq:4.245} y en la ecuación
    \begin{align}
        L''_{xx}(x^k, \lambda^k, \mu^k)d - \sum_{i=1 \atop i \notin I(\bar{x})}^m \mu_i^k g''_i(x^k)d & \nonumber \\
        + (h'(x^k))^\top y + \sum_{i \in I(\bar{x})} z_i^k g'_i(x^k) &= -f'(x^k).
        \label{v2:eq:4.246}
    \end{align}
    Notemos que la única diferencia entre \eqref{v2:eq:4.246} y \eqref{v2:eq:4.232} es el segundo término en el lado izquierdo de \eqref{v2:eq:4.246} (aquí llevamos en cuenta \eqref{v2:eq:4.245}). Mas como $\bar{\mu}_i = 0 \; \forall i \in \{1, \dots, m\} \setminus I(\bar{x})$, existe un número $M > 0$ tal que para todo $x \in \R^n$ suficientemente próximo a $\bar{x}$ y todo $\mu \in \R^m$ suficientemente próximo a $\bar{\mu}$, se tiene que
    \[
        \left\| \sum_{i=1 \atop i \notin I(\bar{x})}^m \mu_i g''_i(x) \right\| \le \sum_{i=1 \atop i \notin I(\bar{x})}^m |\mu_i - \bar{\mu}_i| \|g''_i(x)\| \le M\|\mu - \bar{\mu}\|.
    \]
    Esto significa que el algoritmo cuya iteración es dada por las ecuaciones \eqref{v2:eq:4.232}, \eqref{v2:eq:4.233}, \eqref{v2:eq:4.244} y \eqref{v2:eq:4.245}, puede ser pensado como un método de Newton inexacto con iteraciones determinadas por las ecuaciones \eqref{v2:eq:4.233}, \eqref{v2:eq:4.244}, \eqref{v2:eq:4.245} y \eqref{v2:eq:4.246}. Por el Corolario 3.2.1, propiedades locales y tasa de convergencia de este método son completamente análogas a las del método de Newton.
\end{proof}

Cabe mencionar que la convergencia cuadrática de la sucesión $\{(x^k, \lambda^k, \mu^k)\}$ a $(\bar{x}, \bar{\lambda}, \bar{\mu})$, en principio no implica la convergencia cuadrática (y ni superlineal) de la sucesión $\{x^k\}$ a $\bar{x}$ (vea el Ejemplo 2.2.1). Sin embargo, con frecuencia la convergencia rápida en las variables primales es lo que interesa más en las aplicaciones. Este asunto será investigado a seguir. Al mismo tiempo, consideraremos la posibilidad de escoger las matrices $H_k$ para satisfacer \eqref{v2:eq:4.230} de manera inexacta, siguiendo los principios de la condición de Dennis-Moré (vea el Teorema 3.2.2). Algunas maneras concretas para construir $H_k$ con tales propiedades pueden ser consultadas en [20].

\begin{theorem}[Convergencia superlineal primal del método SQP]\label{v2:thm:4.6.2}
    Supongamos que valen las hipótesis del Teorema 4.6.1 con la excepción, tal vez, de la condición de complementaridad estricta. Supongamos que la sucesión $\{(x^k, \lambda^k, \mu^k)\}$ del Algoritmo 4.6.2 converge a $(\bar{x}, \bar{\lambda}, \bar{\mu})$.
    Entonces la tasa de convergencia de $\{x^k\}$ a $\bar{x}$ es superlineal si, y solamente si,
    \begin{equation}
        P_{K(\bar{x})} \left( (H_k - L''_{xx}(\bar{x}, \bar{\lambda}, \bar{\mu})) (x^{k+1} - x^k) \right) = o(\|x^{k+1} - x^k\|),
        \label{v2:eq:4.247}
    \end{equation}
    donde
    \[
        K(\bar{x}) = \{h \in \ker h'(\bar{x}) \mid \langle g'_i(\bar{x}), h \rangle \le 0 \quad \forall i \in I(\bar{x}), \; \langle f'(\bar{x}), h \rangle \le 0\}
    \]
    es el cono crítico del problema \eqref{v2:eq:4.227}.
\end{theorem}
\begin{proof}
    Recordamos que por la construcción del Algoritmo 4.6.2, el punto $(d^k, \lambda^{k+1}, \mu^{k+1})$ (donde $d^k = x^{k+1} - x^k$) satisface las condiciones KKT para el problema \eqref{v2:eq:4.228}-\eqref{v2:eq:4.229}, i.e.,
    \begin{align}
        f'(x^k) + H_k d^k + (h'(x^k))^\top \lambda^{k+1} + (g'(x^k))^\top \mu^{k+1} &= 0, \label{v2:eq:4.248} \\
        h(x^k) + h'(x^k)d^k &= 0, \label{v2:eq:4.249} \\
        g(x^k) + g'(x^k)d^k &\le 0, \quad \mu^{k+1} \ge 0, \label{v2:eq:4.250} \\
        \mu_i^{k+1}(g_i(x^k) + \langle g'_i(x^k), d^k \rangle) &= 0, \quad i = 1, \dots, m. \label{v2:eq:4.251}
    \end{align}
    De \eqref{v2:eq:4.248} obtenemos que
    \begin{align}
        -H_k d^k &= f'(x^k) + (h'(x^k))^\top \lambda^{k+1} + (g'(x^k))^\top \mu^{k+1} \nonumber \\
        &= L'_x(x^k, \lambda^{k+1}, \mu^{k+1}) \nonumber \\
        &= L'_x(\bar{x}, \lambda^{k+1}, \mu^{k+1}) + L''_{xx}(\bar{x}, \lambda^{k+1}, \mu^{k+1})(x^k - \bar{x}) \nonumber \\
        &\quad + o(\|x^k - \bar{x}\|) - L'_x(\bar{x}, \bar{\lambda}, \bar{\mu}) \nonumber \\
        &= (h'(\bar{x}))^\top(\lambda^{k+1} - \bar{\lambda}) + (g'(\bar{x}))^\top(\mu^{k+1} - \bar{\mu}) \nonumber \\
        &\quad + L''_{xx}(\bar{x}, \bar{\lambda}, \bar{\mu})(x^k - \bar{x}) + o(\|x^k - \bar{x}\|),
        \label{v2:eq:4.252}
    \end{align}
    donde utilizamos las definiciones de punto estacionario del problema \eqref{v2:eq:4.227} y multiplicadores de Lagrange asociados, así como la convergencia de $\{(\lambda^k, \mu^k)\}$ a $(\bar{\lambda}, \bar{\mu})$. De \eqref{v2:eq:4.252}, obtenemos que
    \begin{align}
        (L''_{xx}(\bar{x}, \bar{\lambda}, \bar{\mu}) - H_k)d^k &= (h'(\bar{x}))^\top(\lambda^{k+1} - \bar{\lambda}) + (g'(\bar{x}))^\top(\mu^{k+1} - \bar{\mu}) \nonumber \\
        &\quad + L''_{xx}(\bar{x}, \bar{\lambda}, \bar{\mu})(x^{k+1} - \bar{x}) + o(\|x^k - \bar{x}\|).
        \label{v2:eq:4.253}
    \end{align}
    Supongamos primero que la tasa de convergencia de $\{x^k\}$ a $\bar{x}$ sea superlineal, i.e., $\|x^{k+1} - \bar{x}\| = o(\|x^k - \bar{x}\|)$. En este caso, \eqref{v2:eq:4.253} se escribe como
    \begin{align}
        (L''_{xx}(\bar{x}, \bar{\lambda}, \bar{\mu}) - H_k)d^k &= (h'(\bar{x}))^\top(\lambda^{k+1} - \bar{\lambda}) + (g'(\bar{x}))^\top(\mu^{k+1} - \bar{\mu}) \nonumber \\
        &\quad + o(\|x^k - \bar{x}\|).
        \label{v2:eq:4.254}
    \end{align}
    Como $\{d^k\} \to 0$, $\{\mu^k\} \to \bar{\mu}$ $(k \to \infty)$, de \eqref{v2:eq:4.251} obtenemos que $\mu_i^{k+1} = 0 \; \forall i \in \{1, \dots, m\} \setminus I(\bar{x})$, para todo $k$ suficientemente grande. Utilizando la definición de $\mathcal{K}(\bar{x})$ y la segunda desigualdad en \eqref{v2:eq:4.250}, obtenemos que $\forall v \in \mathcal{K}(\bar{x})$, se tiene que
    \begin{align}
        \langle (h'(\bar{x}))^\top(\lambda^{k+1} - \bar{\lambda}) + (g'(\bar{x}))^\top(\mu^{k+1} - \bar{\mu}), v \rangle \nonumber \\
        = \langle \mu^{k+1}, g'(\bar{x})v \rangle = \sum_{i \in I(\bar{x})} \mu_i^{k+1}\langle g'_i(\bar{x}), v \rangle \le 0.
        \label{v2:eq:4.255}
    \end{align}
    Pelo Teorema 1.3.1(c), la última relación implica en que
    \[
        P_{K(\bar{x})} ((h'(\bar{x}))^\top(\lambda^{k+1} - \bar{\lambda}) + (g'(\bar{x}))^\top(\mu^{k+1} - \bar{\mu})) = 0.
    \]
    De \eqref{v2:eq:4.254} obtenemos ahora que
    \begin{equation}
        P_{K(\bar{x})} ((L''_{xx}(\bar{x}, \bar{\lambda}, \bar{\mu}) - H_k)d^k) = o(\|x^k - \bar{x}\|).
        \label{v2:eq:4.256}
    \end{equation}
    Llevando en cuenta que
    \[
        \|d^k\| \ge \|x^k - \bar{x}\| - \|x^{k+1} - \bar{x}\| = \|x^k - \bar{x}\| + o(\|x^k - \bar{x}\|),
    \]
    la relación \eqref{v2:eq:4.256} implica en \eqref{v2:eq:4.247}, porque
    \[
        \frac{o(\|x^k - \bar{x}\|)}{\|d^k\|} \le \frac{o(\|x^k - \bar{x}\|)}{\|x^k - \bar{x}\| + o(\|x^k - \bar{x}\|)} = \frac{o(\|x^k - \bar{x}\|)/\|x^k - \bar{x}\|}{1 + o(\|x^k - \bar{x}\|)/\|x^k - \bar{x}\|} \to 0 \quad (k \to \infty).
    \]
    Supongamos ahora que valga \eqref{v2:eq:4.247}. Observamos que, por \eqref{v2:eq:4.249}, tem-se que
    \begin{align}
        0 &= h(x^k) + h'(x^k)d^k \nonumber \\
        &= h'(\bar{x})(x^k - \bar{x}) + h'(x^k)d^k + o(\|x^k - \bar{x}\|) \nonumber \\
        &= h'(\bar{x})(x^{k+1} - \bar{x}) + \eta^k,
        \label{v2:eq:4.257}
    \end{align}
    donde, llevando en cuenta que $\{x^k\} \to \bar{x}$,
    \begin{align}
        \eta^k &= (h'(x^k) - h'(\bar{x}))d^k + o(\|x^k - \bar{x}\|) \nonumber \\
        &= o(\|d^k\|) + o(\|x^k - \bar{x}\|).
        \label{v2:eq:4.258}
    \end{align}
    De modo análogo, de la primera desigualdad en \eqref{v2:eq:4.250}, de la relación \eqref{v2:eq:4.251}, y del fato de que $\{\mu^k\} \to \bar{\mu}$, obtenemos que
    \begin{align}
        0 &= \langle g'_i(\bar{x}), x^{k+1} - \bar{x} \rangle + \zeta^k_i \quad \forall i \in I_+(\bar{x}), \label{v2:eq:4.259} \\
        0 &\ge \langle g'_i(\bar{x}), x^{k+1} - \bar{x} \rangle + \zeta^k_i \quad \forall i \in I_0(\bar{x}), \label{v2:eq:4.260} \\
        0 &= \mu_i^{k+1}(\langle g'_i(\bar{x}), x^{k+1} - \bar{x} \rangle + \zeta^k_i) \quad \forall i \in I_0(\bar{x}), \label{v2:eq:4.261}
    \end{align}
    donde $I_+(\bar{x}) = \{i \in I(\bar{x}) \mid \bar{\mu}_i > 0\}$, $I_0(\bar{x}) = \{i \in I(\bar{x}) \mid \bar{\mu}_i = 0\}$,
    \begin{equation}
        \zeta^k_i = o(\|d^k\|) + o(\|x^k - \bar{x}\|), \quad i \in I(\bar{x}).
        \label{v2:eq:4.262}
    \end{equation}
    Por la condición de independencia lineal de los gradientes de las restricciones activas \eqref{v1:eq:1.17} y por las relaciones \eqref{v2:eq:4.258} y \eqref{v2:eq:4.262}, existe un elemento $\xi^k \in \R^n$ tal que
    \begin{align}
        h'(\bar{x})\xi^k &= \eta^k, \quad \langle g'_i(\bar{x}), \xi^k \rangle = \zeta^k_i \quad \forall i \in I(\bar{x}), \label{v2:eq:4.263} \\
        \|\xi^k\| &= o(\|d^k\|) + o(\|x^k - \bar{x}\|). \label{v2:eq:4.264}
    \end{align}
    Definimos $v^k = x^{k+1} - \bar{x} + \xi^k$. De \eqref{v2:eq:4.257}-\eqref{v2:eq:4.260} y de \eqref{v2:eq:4.263}, tenemos que $v^k \in \mathcal{K}(\bar{x})$. En particular, de modo análogo a \eqref{v2:eq:4.255}, obtenemos que
    \begin{align*}
        &\langle (h'(\bar{x}))^\top(\lambda^{k+1} - \bar{\lambda}) + (g'(\bar{x}))^\top(\mu^{k+1} - \bar{\mu}), v^k \rangle \\
        &= \sum_{i \in I(\bar{x})} \mu_i^{k+1}\langle g'_i(\bar{x}), v^k \rangle = 0,
    \end{align*}
    donde la última igualdad sigue de \eqref{v2:eq:4.259} y \eqref{v2:eq:4.261}. Multiplicando escalarmente los dos lados de \eqref{v2:eq:4.253} por $v^k$ y utilizando la relación de arriba, obtenemos
    \begin{equation}
        \langle (L''_{xx}(\bar{x}, \bar{\lambda}, \bar{\mu}) - H_k)d^k, v^k \rangle = \langle L''_{xx}(\bar{x}, \bar{\lambda}, \bar{\mu})(x^{k+1} - \bar{x}), v^k \rangle + o(\|x^k - \bar{x}\|\|v^k\|).
        \label{v2:eq:4.265}
    \end{equation}
    Pela condición suficiente de segunda orden existe un número $\gamma > 0$ tal que
    \[
        \langle L''_{xx}(\bar{x}, \bar{\lambda}, \bar{\mu})v, v \rangle \ge \gamma \|v\|^2 \quad \forall v \in \mathcal{K}(\bar{x}).
    \]
    Utilizando también \eqref{v2:eq:4.265}, obtenemos
    \begin{align*}
        \gamma\|v^k\|^2 &\le \langle L''_{xx}(\bar{x}, \bar{\lambda}, \bar{\mu})v^k, v^k \rangle \\
        &= \langle (L''_{xx}(\bar{x}, \bar{\lambda}, \bar{\mu}) - H_k)d^k, v^k \rangle + \langle L''_{xx}(\bar{x}, \bar{\lambda}, \bar{\mu})\xi^k, v^k \rangle \\
        &\quad + o(\|x^k - \bar{x}\|\|v^k\|) \\
        &= \langle (L''_{xx}(\bar{x}, \bar{\lambda}, \bar{\mu}) - H_k)d^k, v^k \rangle + o(\|d^k\|\|v^k\|) + o(\|x^k - \bar{x}\|\|v^k\|) \\
        &\le \langle P_{\mathcal{K}(\bar{x})} ((L''_{xx}(\bar{x}, \bar{\lambda}, \bar{\mu}) - H_k)d^k), v^k \rangle \\
        &\quad + o(\|d^k\|\|v^k\|) + o(\|x^k - \bar{x}\|\|v^k\|) \\
        &= o(\|d^k\|\|v^k\|) + o(\|x^k - \bar{x}\|\|v^k\|),
    \end{align*}
    donde la segunda igualdad sigue de \eqref{v2:eq:4.264}, la segunda desigualdad sigue del Teorema 1.3.1(c), y la última igualdad sigue de \eqref{v2:eq:4.247}. Concluimos que $\|v^k\| = o(\|d^k\|) + o(\|x^k - \bar{x}\|)$. Mas en este caso, por \eqref{v2:eq:4.264}, obtenemos que
    \[
        \|x^{k+1} - \bar{x}\| = \|v^k - \xi^k\| = o(\|d^k\|) + o(\|x^k - \bar{x}\|).
    \]
    Esto significa que existen sucesiones $\{t_k\} \subset \R_+$ y $\{s_k\} \subset \R_+$ tales que $t_k \to 0, s_k \to 0$ $(k \to \infty)$ y, para todo $k$,
    \begin{align*}
        \|x^{k+1} - \bar{x}\| &\le t_k\|x^{k+1} - x^k\| + s_k\|x^k - \bar{x}\| \\
        &\le \max\{t_k, s_k\} (\|x^{k+1} - \bar{x}\| + 2\|x^k - \bar{x}\|).
    \end{align*}
    De aquí, para todo $k$ suficientemente grande,
    \[
        \|x^{k+1} - \bar{x}\| \le \frac{2\max\{t_k, s_k\}}{1 - \max\{t_k, s_k\}} \|x^k - \bar{x}\|,
    \]
    y como $\max\{t_k, s_k\} \to 0$ $(k \to \infty)$, tenemos que
    \[
        x^{k+1} - \bar{x} = o(\|x^k - \bar{x}\|),
    \]
    i.e., la convergencia de $\{x^k\}$ a $\bar{x}$ es superlineal.
\end{proof}

Observamos que en el caso de la elección de la matriz $H_k$ de acuerdo con \eqref{v2:eq:4.230}, la convergencia en las variables primales ciertamente es superlineal, ya que $H_k \to L''_{xx}(\bar{x}, \bar{\lambda}, \bar{\mu})$ cuando $(x^k, \lambda^k, \mu^k) \to (\bar{x}, \bar{\lambda}, \bar{\mu})$ $(k \to \infty)$ y, por tanto, la validez de la condición \eqref{v2:eq:4.247} es obvia.
Las hipótesis del Teorema 4.6.1 que garantizan convergencia local superlineal (en variables primales-duales) del método SQP pueden ser relajadas. Por ejemplo, es posible eliminar la condición de complementaridad estricta, mas a costa de la introducción de una condición de segunda orden más fuerte. Esta condición, llamada de \textit{condición suficiente de segunda orden fuerte}, consiste en lo siguiente:
\[
    \langle L''_{xx}(\bar{x}, \bar{\lambda}, \bar{\mu})v, v \rangle > 0 \quad \forall v \in \mathcal{K}_+(\bar{x}) \setminus \{0\},
\]
donde
\[
    \mathcal{K}_+(\bar{x}) = \{v \in \ker h'(\bar{x}) \mid \langle g'_i(\bar{x}), v \rangle = 0 \quad \forall i \in I_+(\bar{x})\};
\]
\[
    I_+(\bar{x}) = \{i \in I(\bar{x}) \mid \bar{\mu}_i > 0\}.
\]
En el caso de complementaridad estricta, tem-se que $I_+(\bar{x}) = I(\bar{x})$ y, por tanto, $\mathcal{K}_+(\bar{x}) = \mathcal{K}(\bar{x})$ y la condición suficiente de segunda orden fuerte coincide con la condición suficiente de segunda orden común.
Investigaciones contemporáneas del método SQP y de diversas modificaciones de él tienen por objetivo la relajación de las hipótesis utilizadas para garantizar convergencia local superlineal. Desenvolvimientos en esta dirección hacen uso de técnicas bien más sofisticadas, que están fuera del alcance de este libro. El resultado más fuerte sobre convergencia local superlineal del método SQP (conocido actualmente) está basado en la condición suficiente de segunda orden (común) y en la hipótesis de que el multiplicador de Lagrange asociado a la solución sea único (esta última condición también se llama condición de Mangasarian-Fromovitz estricta; vea [12]).
Finalmente, comentaremos sobre la comparación de los métodos de Newton generalizados para el sistema KKT en el \S 4.5 con el método SQP. La convergencia local cuadrática de los métodos de Newton generalizados fue establecida bajo las hipótesis de independencia lineal de los gradientes de las restricciones activas y la condición suficiente de segunda orden fuerte (vea el Teorema 4.5.2), en cuanto el método SQP requiere la unicidad de multiplicador de Lagrange y la condición suficiente de segunda orden usual, que son hipótesis bien más débiles. En compensación, una iteración del método de Newton generalizado consiste en la resolución de un sistema de ecuaciones lineales, lo que tiene un costo computacional bien menor que la resolución de un problema de programación cuadrática en SQP. Notemos también que la condición (4.161) en el Teorema 4.5.3, que es necesaria para la convergencia superlineal en las variables primales del método de Newton generalizado, es un poco más débil que la condición correspondiente en el caso de SQP (vea el Teorema 4.6.2). La estructura de SQP permite varias estrategias de globalización bien naturales (que serán presentadas a seguir). La globalización del método de Newton generalizado también es posible, mas ella es más difícil y generalmente es menos eficiente desde el punto de vista computacional.
%%% Local Variables:
%%% mode: LaTeX
%%% TeX-master: "../../../Apuntes-Programacion-No-Lineal"
%%% End:
